% 05-ket-qua.tex — Chương 5: Kết quả và đánh giá
\chapter{Kết quả và đánh giá}
\label{ch:ketqua}

Chương này trình bày kết quả triển khai theo hướng định tính: mô tả những gì đã thực hiện
được và các nhận định rút ra từ quan sát thực tế. Mục \ref{sec:pp-danhgia} trình bày phương
pháp đánh giá đã thiết kế; mục \ref{sec:kq-retrieval} trình bày kết quả khâu retrieval; mục
\ref{sec:kq-generation} trình bày khâu sinh văn bản; mục \ref{sec:ss} so sánh các hướng
tiếp cận; mục \ref{sec:hieunang} trình bày hiệu năng hệ thống; mục \ref{sec:casestudy} nêu
case study; mục \ref{sec:hieuqua} trình bày hiệu quả vận hành; và mục \ref{sec:thaoluan}
thảo luận.

\begin{ghichu}
Ghi chú minh bạch: tại thời điểm viết báo cáo, việc tổng hợp các chỉ số định lượng trên
toàn bộ bộ đánh giá chưa hoàn tất. Do đó chương này chỉ trình bày kết quả định tính và
những việc đã làm được; các chỉ số chi tiết sẽ được công bố ở đợt đánh giá tiếp theo.
\end{ghichu}

\section{Phương pháp đánh giá}
\label{sec:pp-danhgia}

Đề tài thiết kế đánh giá theo ba lớp:

\textbf{Lớp retrieval:} quy trình tái lập chính xác đường đi sản xuất -- với mỗi câu hỏi,
hệ thống chạy sparse retrieval BM25 và dense retrieval, lấy hợp ứng viên, rerank, rồi so
sánh tập kết quả với đáp án tham chiếu. Các thước đo được sử dụng gồm Recall@$k$ cho từng
nhánh, Recall sau khi hợp, Recall@$k$ sau rerank và MRR. Điểm nổi bật của quy trình là
việc \emph{phân tách lỗi} thành hai loại: lỗi \emph{RECALL} (chunk đúng không được truy
xuất) và lỗi \emph{RERANK} (chunk đúng có trong ứng viên nhưng bị rerank loại khỏi
top-$k$). Cách phân tách này giúp định vị điểm nghẽn chính xác thay vì chỉ nhìn vào một
con số tổng.

\textbf{Lớp sinh văn bản:} hai phương pháp. Thứ nhất, LLM-as-judge cho điểm câu trả lời
theo thang điểm, dựa trên so sánh với đáp án tham chiếu. Thứ hai, bộ thước đo RAGAS gồm
bốn thành phần (Faithfulness, Answer Relevancy, Context Precision, Context Recall), được
cài đặt theo hai phiên bản: một bản dùng thư viện RAGAS chính thức và một bản tự cài đặt
để kiểm soát chi tiết từng thành phần, với mô hình đánh giá độc lập.

\textbf{Lớp hiệu năng:} đo throughput và latency của vLLM ở hai chế độ -- một yêu cầu đơn
(single-stream) và lô lớn (batch) -- nhằm xác nhận hạ tầng đáp ứng được nhu cầu phục vụ
đồng thời.

\section{Kết quả khâu retrieval}
\label{sec:kq-retrieval}

\subsection{Đã thực hiện được}

Bộ đánh giá retrieval đã được xây dựng hoàn chỉnh: có dataset với đáp án tham chiếu, quy
trình chạy đúng như cấu hình sản xuất, và xuất kết quả ra file để phân tích. Quy trình
này chạy được độc lập với API, giúp kiểm tra lại mọi thay đổi cấu hình (số ứng viên mỗi
nhánh, số chunk sau rerank, ngưỡng điểm) mà không ảnh hưởng hệ thống đang chạy.

\subsection{Nhận định rút ra}

Qua quan sát kết quả, một số nhận định đáng chú ý:

\begin{itemize}
  \item \textbf{Dense retrieval mạnh ở diễn đạt khác từ khóa.} Với những câu hỏi diễn đạt
  khác với cách viết trong tài liệu (người dùng hỏi theo cách hiểu của mình), dense
  retrieval truy xuất được chunk phù hợp trong khi BM25 bỏ sót. Đây là ưu thế cốt lõi
  khiến dense retrieval trở thành nhánh chính.
  \item \textbf{BM25 mạnh với từ khóa chính xác và số hiệu văn bản.} Khi người dùng nêu
  đúng số hiệu quyết định, mã dịch vụ hoặc thuật ngữ hiếm, BM25 khớp chính xác hơn.
  Vì vậy nhánh BM25 vẫn có vai trò bổ trợ, không thể loại bỏ.
  \item \textbf{Hybrid retrieval đảm bảo phạm vi ứng viên.} Việc hợp hai nhánh giúp tập
  ứng viên chứa được chunk đúng ở cả hai loại câu hỏi, là điều kiện để bước sau có cơ hội
  xếp đúng.
  \item \textbf{Rerank là điểm cần cải thiện.} Khi số chunk sau rerank nhỏ, một phần chunk
  đúng bị loại khỏi ngữ cảnh. Phân tách lỗi cho thấy đây là lỗi thuộc loại \emph{RERANK},
  không phải \emph{RECALL}: chunk đúng luôn nằm trong tập ứng viên, nhưng bước xếp hạng
  chưa đưa nó vào top-$k$. Nguyên nhân khả dĩ là reranker ưu tiên các chunk ``trông
  giống'' câu hỏi về mặt từ vựng nhưng thiếu thông tin đáp án.
\end{itemize}

Nhận định cuối cùng có ý nghĩa thực tiễn trực tiếp: thay vì cố cải thiện retrieval, nỗ lực
cải tiến nên tập trung vào reranker -- ví dụ fine-tune trên dữ liệu tiếng Việt miền bưu
chính, điều chỉnh số chunk sau rerank, hoặc kết hợp điểm rerank với điểm truy xuất ban
đầu. Đây chính là nội dung thiết kế fine-tuning reranker ở chương \ref{ch:kientruc}.

\section{Kết quả khâu sinh văn bản}
\label{sec:kq-generation}

\subsection{Đã thực hiện được}

Hai bộ công cụ đánh giá sinh văn bản đã được cài đặt và chạy được: LLM-as-judge cho điểm
câu trả lời theo thang điểm dựa trên đáp án tham chiếu, và bộ thước đo RAGAS cho các
thành phần Faithfulness, Answer Relevancy, Context Precision, Context Recall. Bộ đánh giá
cho miền EMS gồm 99 câu hỏi -- đáp án tham chiếu đã được chuẩn bị, và hệ thống đã thu thập
đầy đủ phản hồi cho toàn bộ 99 câu: mỗi bản ghi gồm câu hỏi, đáp án tham chiếu, câu trả
lời của chatbot, ngữ cảnh retrieval và độ dài ký tự.

\subsection{Nhận định rút ra}

Qua quan sát trực tiếp các hội thoại, hệ thống cho thấy: câu trả lời bám sát tài liệu,
trích dẫn được nguồn (tiêu đề/số hiệu văn bản), và \emph{biết từ chối} khi không tìm thấy
bằng chứng phù hợp thay vì đoán. Hành vi từ chối này có ý nghĩa quan trọng đối với doanh
nghiệp, vì câu trả lời sai về biểu phí hoặc chính sách bồi thường dễ gây khiếu nại.

Việc tổng hợp các điểm RAGAS từ phản hồi đã thu thập chưa hoàn tất tại thời điểm viết
báo cáo, do đó báo cáo chưa công bố các giá trị tổng hợp. Đây là việc sẽ được thực hiện
trong đợt đánh giá tiếp theo, cùng với việc mở rộng bộ câu hỏi.

\section{So sánh các hướng tiếp cận}
\label{sec:ss}

Bảng \ref{tab:ss} so sánh ba hướng tiếp cận dựa trên đặc điểm triển khai và quan sát thực
tế. Với \rag, các nhận định dựa trên kết quả triển khai; với \ft{} và kiến trúc hybrid,
các nhận định mang tính phân tích thiết kế do chưa có kết quả huấn luyện thật.

\begin{table}[htbp]
\centering
\caption{So sánh RAG, fine-tuning và kiến trúc hybrid trong bối cảnh VNPost.}
\label{tab:ss}
\small
\begin{tabular}{@{}p{3.4cm}p{3.9cm}p{3.6cm}p{3.6cm}@{}}
\toprule
\textbf{Tiêu chí} & \textbf{RAG} & \textbf{Fine-tuning} & \textbf{Hybrid} \\
\midrule
Tri thức động & Rất tốt & Kém & Tốt \\
Trích nguồn & Có & Hạn chế & Có \\
Phong cách miền & Trung bình & Tốt & Tốt \\
Chi phí cập nhật & Thấp & Cao & Thấp--trung bình \\
Độ trễ & Cao hơn & Thấp & Trung bình \\
Độ phức tạp & Trung bình & Cao & Cao \\
Trạng thái đánh giá & Đã triển khai & Đề xuất & Đề xuất \\
\bottomrule
\end{tabular}
\end{table}

Kết luận chính: \rag{} là lựa chọn phù hợp và khả thi nhất cho nhu cầu hiện tại của VNPost
vì tri thức thay đổi thường xuyên và cần trích nguồn; \ft{} bổ trợ ở lớp hành vi; và kiến
trúc hybrid là đích đến dài hạn, khi dữ liệu hội thoại đã đủ chín để huấn luyện.

\section{Hiệu năng hệ thống}
\label{sec:hieunang}

Hạ tầng phục vụ được xây dựng quanh vLLM với cơ chế PagedAttention và quantization FP8,
cho phép phục vụ các mô hình cỡ 20--35 tỷ tham số trên GPU của hệ thống nội bộ. Việc đo
throughput ở chế độ đơn và lô đã được thiết lập, với kết quả cho thấy mức thông lượng đủ
để phục vụ đồng thời nhu cầu của một tổng công ty, và latency chấp nhận được cho câu trả
lời ngắn trong vài giây.

Điểm đáng lưu ý là quantization FP8 giảm đáng kể bộ nhớ và tăng tốc tính toán với mức suy
giảm chất lượng nhỏ, giúp số lượng mô hình phục vụ song song (LLM, embedding, reranker,
OCR) nằm trong khả năng của hạ tầng. Kết quả đo chi tiết các cấu hình sẽ được báo cáo khi
hoàn tất chuẩn hóa quy trình benchmark.

\section{Case study hội thoại}
\label{sec:casestudy}

Để minh họa hành vi thực tế, dưới đây là một số tình huống hội thoại ví dụ (không phải bản
ghi hội thoại thật).

\paragraph{Tình huống 1 -- Hỏi biểu phí.} Người dùng hỏi cước chuyển phát cho một trọng
lượng và quốc gia cụ thể. Agent điều phối định tuyến đến agent \rag{} sản phẩm -- dịch vụ;
hệ thống truy xuất các chunk biểu phí liên quan, rerank, rồi sinh câu trả lời kèm trích
dẫn số hiệu quyết định. Nếu không tìm thấy ngữ cảnh đủ tin cậy, hệ thống thông báo thiếu
thông tin thay vì đoán.

\paragraph{Tình huống 2 -- Tra cứu bưu gửi.} Người dùng cung cấp mã bưu gửi. Agent điều
phối gọi agent tra cứu bưu gửi, agent này gọi công cụ tracking qua SOAP/API và trả về tình
trạng mới nhất cùng các mốc thời gian.

\paragraph{Tình huống 3 -- Tra cứu bưu cục.} Người dùng hỏi địa chỉ bưu cục gần một địa
điểm. Agent tra cứu bưu cục truy vấn cơ sở dữ liệu với tìm kiếm mờ để xử lý sai chính tả
và trả về danh sách bưu cục phù hợp.

\paragraph{Tình huống 4 -- Chuyển hỗ trợ thủ công.} Khi yêu cầu vượt ngoài phạm vi hoặc
cần xử lý nghiệp vụ phức tạp, agent điều phối tạo ticket và chuyển đến nhân viên, bảo đảm
không bỏ sót yêu cầu.

\section{Hiệu quả vận hành}
\label{sec:hieuqua}

Về mặt vận hành, giải pháp mang lại hiệu quả rõ rệt ở các khía cạnh sau:

\begin{itemize}
  \item \textbf{Giảm thời gian tra cứu:} cán bộ nghiệp vụ không phải mở nhiều văn bản để
  tìm điều khoản; câu trả lời kèm nguồn được trả về ngay trong hội thoại.
  \item \textbf{Chuẩn hóa câu trả lời:} cùng một câu hỏi nhận được câu trả lời nhất quán,
  theo đúng văn bản hiệu lực, giảm sai lệch giữa các điểm giao dịch.
  \item \textbf{Giảm tải kênh hỗ trợ thủ công:} phần câu hỏi tra cứu đơn giản được giải
  quyết bởi hệ thống, để nhân lực tập trung vào các yêu cầu phức tạp.
  \item \textbf{Hỗ trợ dữ liệu và báo cáo:} kết hợp với các agent tra cứu chất lượng, hệ
  thống giúp lập báo cáo và phân tích nhanh, hỗ trợ quyết định vận hành.
\end{itemize}

Ước tính định lượng về lợi ích kinh tế -- xã hội (thời gian tiết kiệm, quy mô người dùng,
giá trị quy đổi) sẽ được thực hiện khi có dữ liệu vận hành thực tế sau giai đoạn triển
khai rộng, nhằm bảo đảm con số có căn cứ và có thể kiểm chứng.

\section{Thảo luận}
\label{sec:thaoluan}

\subsection{Ưu điểm}

Hệ thống cho thấy RAG hybrid với rerank là giải pháp khả thi cho tri thức nội bộ động của
VNPost: dense retrieval mạnh về ngữ nghĩa, BM25 bổ trợ cho từ khóa chính xác, kiến trúc
multi-agent giúp mở rộng theo nghiệp vụ, và việc phục vụ cục bộ đáp ứng yêu cầu bảo mật.
Quan sát tập trung và quản lý prompt giúp vận hành và cải tiến liên tục.

\subsection{Hạn chế}

Thứ nhất, việc tổng hợp chỉ số định lượng trên toàn bộ bộ đánh giá chưa hoàn tất, nên các
kết luận chủ yếu mang tính định tính. Thứ hai, khâu rerank còn hạn chế và làm mất một phần
chunk đúng ở top-$k$ nhỏ. Thứ ba, phần \ft{} chưa được huấn luyện thật nên chưa có bằng
chứng định lượng. Thứ tư, đánh giá chủ yếu dựa trên mô hình (LLM-as-judge, RAGAS) mà chưa
có đánh giá diện rộng với người dùng cuối.

\subsection{Tính khái quát}

Kiến trúc và quy trình đánh giá đề xuất có thể áp dụng cho các tổ chức có đặc thù tương
tự: khối tài liệu nội bộ lớn, nhu cầu tra cứu thường xuyên, và yêu cầu bảo mật dữ liệu.
Việc thay thế kho tài liệu và agent chuyên biệt cho phép tái sử dụng nền tảng cho miền
khác.

\section*{Kết chương}
\addcontentsline{toc}{section}{Kết chương}
Chương này đã trình bày những việc đã làm được ở khâu retrieval, sinh văn bản, hiệu năng và
case study, cùng các nhận định định tính. Chương tiếp theo kết luận và đề xuất hướng phát
triển.
